\documentclass[11pt,a4paper]{article}

% ---- idioma y encoding ----
\usepackage[utf8]{inputenc}
\usepackage[spanish]{babel}
\usepackage[T1]{fontenc}

% ---- margenes y layout ----
\usepackage{amsmath,amssymb}
\usepackage[margin=2.2cm]{geometry}
\usepackage{setspace}
\onehalfspacing

% ---- colores ----
\usepackage[dvipsnames]{xcolor}
\definecolor{codebg}{HTML}{F5F5FA}
\definecolor{codeframe}{HTML}{CCCCCC}
\definecolor{accentblue}{HTML}{3B82F6}
\definecolor{darktext}{HTML}{1E1E1E}
\definecolor{graytext}{HTML}{666666}
\definecolor{tabheader}{HTML}{E8EDF5}
\definecolor{lightgreen}{HTML}{D4EDDA}
\definecolor{lightyellow}{HTML}{FFF3CD}

% ---- listings para codigo ----
\usepackage{listings}
\lstset{
    basicstyle=\ttfamily\small,
    backgroundcolor=\color{codebg},
    frame=single,
    rulecolor=\color{codeframe},
    breaklines=true,
    breakatwhitespace=true,
    tabsize=4,
    showstringspaces=false,
    keywordstyle=\color{accentblue}\bfseries,
    commentstyle=\color{graytext}\itshape,
    stringstyle=\color{OliveGreen},
    numbers=none,
    xleftmargin=8pt,
    xrightmargin=8pt,
    aboveskip=6pt,
    belowskip=6pt,
    framexleftmargin=4pt,
    framextopmargin=3pt,
    framexbottommargin=3pt,
    lineskip=-1pt,
    columns=fullflexible,
}

% ---- tablas ----
\usepackage{booktabs}
\usepackage{array}
\usepackage{tabularx}
\usepackage{colortbl}

% ---- headers/footers ----
\usepackage{fancyhdr}
\pagestyle{fancy}
\fancyhf{}
\fancyhead[L]{\small\textcolor{graytext}{Guia Completa — INFO1167 Robotica}}
\fancyhead[R]{\small\textcolor{graytext}{Proyectos 2 y 3}}
\fancyfoot[C]{\thepage}
\renewcommand{\headrulewidth}{0.4pt}

% ---- hipervinculos ----
\usepackage[colorlinks=true,linkcolor=accentblue,urlcolor=accentblue]{hyperref}

% ---- boxes ----
\usepackage{tcolorbox}
\tcbuselibrary{skins,breakable}

\newtcolorbox{importante}{
    colback=lightyellow,
    colframe=OliveGreen,
    fonttitle=\bfseries,
    title=Punto Clave,
    breakable,
    boxrule=0.8pt,
    arc=3pt,
    left=8pt, right=8pt, top=6pt, bottom=6pt
}

\newtcolorbox{definicion}[1][]{
    colback=blue!5,
    colframe=accentblue,
    fonttitle=\bfseries,
    title={#1},
    breakable,
    boxrule=0.8pt,
    arc=3pt,
    left=8pt, right=8pt, top=6pt, bottom=6pt
}

\newtcolorbox{pregunta}[1][]{
    colback=green!5,
    colframe=OliveGreen,
    fonttitle=\bfseries,
    title={#1},
    breakable,
    boxrule=0.8pt,
    arc=3pt,
    left=8pt, right=8pt, top=6pt, bottom=6pt
}

% ---- comandos utiles ----
\newcommand{\codigo}[1]{\texttt{\small #1}}
\newcommand{\vocab}[1]{\textbf{\textcolor{accentblue}{#1}}}

% ====================================================================
\begin{document}

% ---- PORTADA ----
\begin{titlepage}
\centering
\vspace*{3cm}
{\Huge\bfseries\textcolor{darktext}{Guia Completa}\par}
\vspace{0.8cm}
{\LARGE\textcolor{graytext}{Proyectos 2 y 3 de Robotica}\par}
\vspace{0.5cm}
{\Large\textcolor{graytext}{INFO1167 — Alberto Caro}\par}
\vspace{2cm}
{\large\textcolor{graytext}{Todo lo que necesitas saber para defender en 2 dias}\par}
\vspace{0.3cm}
{\large\textcolor{graytext}{Explicado como si no supieras nada}\par}
\vfill
{\large\today\par}
\end{titlepage}

% ---- TABLA DE CONTENIDOS ----
\tableofcontents
\newpage

% ====================================================================
\section{Conceptos Basicos}
% ====================================================================

\subsection{Que es una grilla?}

Imagina un tablero de ajedrez pero mas chico: \textbf{6 filas por 7 columnas}. Cada casilla puede ser:

\begin{itemize}
    \item \textbf{Libre (S)} — el robot puede caminar ahi
    \item \textbf{Muro (X)} — obstaculo, no puede pasar
    \item \textbf{Meta (M)} — el objetivo, donde tiene que llegar
    \item \textbf{Vacio (.)} — no existe, no es parte del mapa
\end{itemize}

Nuestro mapa se ve asi:

\begin{lstlisting}
      0  1  2  3  4  5  6
  0   .  S  S  S  S  .  .
  1   .  S  .  .  S  .  .
  2   S  S  S  X  S  S  S
  3   .  .  S  .  S  .  .
  4   .  .  S  S  S  .  .
  5   .  .  .  M  .  .  .
\end{lstlisting}

El robot empieza en una posicion \textbf{random} de las casillas \codigo{S} y tiene que llegar a \codigo{M}. Hay 18 estados transitables y 1 muro.

\subsection{Conceptos clave}

\begin{definicion}[Estado]
Un \vocab{estado} es simplemente ``donde esta el robot ahora''. Si esta en fila 2, columna 1, su estado es \codigo{(2, 1)}. Cada casilla libre es un estado posible.
\end{definicion}

\begin{definicion}[Accion]
Una \vocab{accion} es un movimiento: Norte, Sur, Este u Oeste. El robot elige una accion en cada estado.
\end{definicion}

\begin{definicion}[Politica]
La \vocab{politica} es la lista de ``si estoy aca, voy para alla''. Para cada casilla, dice cual es el mejor movimiento. Como un GPS que siempre te dice para donde doblar.
\end{definicion}

\begin{definicion}[Algoritmo]
Un \vocab{algoritmo} es una receta paso a paso. ``Primero hace esto, despues hace esto otro, y si pasa esto, hace esto otro.''
\end{definicion}

\newpage

% ====================================================================
\section{Proyecto 2 — MDP y SMDP}
% ====================================================================

\subsection{Que es un MDP?}

\textbf{MDP} = \textit{Markov Decision Process} (Proceso de Decision de Markov).

Imagina que estas en un laberinto oscuro. No puedes ver nada, pero tienes un mapa mental. Cada vez que das un paso:

\begin{itemize}
    \item Pierdes \textbf{1 punto} (costo de caminar)
    \item Si llegas a la salida, ganas \textbf{10 puntos}
    \item A veces resbalas y no te mueves (\textbf{10\%} de las veces)
\end{itemize}

El MDP te dice: ``desde cada casilla, ¿para donde debo ir para ganar la mayor cantidad de puntos?''

\subsubsection{Los 5 ingredientes del MDP}

\begin{center}
\begin{tabular}{>{\bfseries}l l l}
\toprule
\rowcolor{tabheader}
Ingrediente & Significado & En nuestro juego \\
\midrule
S (Estados) & Todos los lugares posibles & Las 18 casillas libres \\
A (Acciones) & Lo que puede hacer & Norte, Sur, Este, Oeste \\
P(s'$|$s,a) & Probabilidad de transicion & 90\% exito, 10\% falla \\
R(s,a) & Recompensa & $-1$ por paso, $+10$ en meta \\
$\gamma$ & Factor de descuento & 0.97 \\
\bottomrule
\end{tabular}
\end{center}

\subsection{La ecuacion de Bellman}

Esta es la ecuacion \textbf{mas importante} del MDP:

\begin{equation}
V(s) = \max_a \left[ R(s,a) + \gamma \sum_{s'} P(s'|s,a) \cdot V(s') \right]
\end{equation}

Desglose termino por termino:

\begin{importante}
\begin{itemize}
    \item $V(s)$ = ``que tan bueno es estar en esta casilla''
    \item $\max_a$ = ``elijo la mejor accion de todas''
    \item $R(s,a)$ = ``lo que gano/pierdo ahora'' ($-1$ por caminar)
    \item $\gamma = 0.97$ = ``cuanto me importa el futuro''
    \item $V(s')$ = ``que tan bueno es el lugar donde voy a llegar''
\end{itemize}
\end{importante}

La ecuacion dice: \textbf{``el valor de donde estoy = lo que gano ahora + lo que espero ganar en el futuro, eligiendo la mejor accion''}

\subsection{Factor de descuento $\gamma = 0.97$}

¿Por que 0.97 y no 1.0? Porque si $\gamma = 1.0$, el robot valdria igual una recompensa hoy que en 1000 pasos. Con $\gamma = 0.97$:

\begin{center}
\begin{tabular}{l l}
\toprule
\rowcolor{tabheader}
Momento & Valor \\
\midrule
Ahora & $1.0$ \\
En 1 paso & $0.97$ \\
En 2 pasos & $0.97^2 = 0.94$ \\
En 10 pasos & $0.97^{10} = 0.74$ \\
\bottomrule
\end{tabular}
\end{center}

El robot prefiere recompensas cercanas pero no ignora las lejanas.

\subsection{Probabilidad de exito (90\%)}

Cuando el robot decide ir al Sur:
\begin{itemize}
    \item \textbf{90\%} de las veces: se mueve al Sur correctamente
    \item \textbf{10\%} de las veces: se queda donde esta (fallo)
\end{itemize}

\begin{lstlisting}[language=Python]
valor = 0.9 * (recompensa + gamma * V[siguiente])   # exito
      + 0.1 * (recompensa + gamma * V[mismo_lugar])  # fallo
\end{lstlisting}

\newpage
\subsection{El codigo del MDP}

La funcion clave del proyecto 2:

\begin{lstlisting}[language=Python]
def calcular_politica(gamma, modo='MDP'):
    V = {s: 0.0 for s in estados}   # todos empiezan en 0
    for _ in range(500):             # iterar hasta converger
        for s in estados:
            if s == meta:
                V[meta] = 10.0       # la meta vale 10
                continue
            mejor = float('-inf')
            for acc in ACTIONS:       # probar las 4 direcciones
                sig = move(s, acc)    # donde llegaria
                r = -1.0 + (10.0 if sig == meta else 0.0)

                # AQUI esta la diferencia MDP vs SMDP
                if modo == 'SMDP':
                    desc = promedio(gamma ** tau)   # temporal
                else:
                    desc = gamma                    # fijo

                val = 0.9*(r + desc*V[sig]) + 0.1*(r + desc*V[s])
                mejor = max(mejor, val)
            V[s] = mejor

    # extraer politica: mejor accion para cada estado
    for s in estados:
        politica[s] = accion_con_mayor_valor
\end{lstlisting}

Paso a paso:
\begin{enumerate}
    \item Empezamos con $V(s) = 0$ para todas las casillas
    \item Para cada casilla, probamos las 4 acciones
    \item Para cada accion, calculamos cuanto vale usando Bellman
    \item Nos quedamos con la mejor accion
    \item Repetimos hasta que los valores se estabilicen (convergencia)
    \item Al final, la politica es: para cada casilla, la accion con mayor valor
\end{enumerate}

\newpage

\subsection{Que es un SMDP?}

\textbf{SMDP} = \textit{Semi-Markov Decision Process}.

La unica diferencia con el MDP es que las acciones toman \textbf{tiempo variable}:

\begin{center}
\begin{tabular}{l l}
\toprule
\rowcolor{tabheader}
Accion & Tiempo \\
\midrule
Norte & Normal($\mu=2, \sigma=0.2$) \\
Sur & Normal($\mu=2, \sigma=0.2$) \\
Este & Normal($\mu=3, \sigma=0.3$) \\
Oeste & Normal($\mu=3, \sigma=0.3$) \\
\bottomrule
\end{tabular}
\end{center}

\subsubsection{Por que importa el tiempo?}

Porque el factor de descuento penaliza el tiempo. Si una accion toma mas tiempo:

\begin{lstlisting}
gamma^tau = 0.97^3 = 0.91   (accion lenta, descuenta mucho)
gamma^tau = 0.97^2 = 0.94   (accion rapida, descuenta poco)
\end{lstlisting}

\begin{importante}
El SMDP prefiere acciones rapidas cuando ambas llevan a buen lugar. Norte/Sur son mas rapidos (2s) que Este/Oeste (3s), entonces si ambas opciones llegan a la meta, el SMDP prefiere Norte/Sur.
\end{importante}

\subsubsection{Diferencia en el codigo}

\begin{lstlisting}[language=Python]
# MDP: descuento fijo
desc = gamma                    # siempre 0.97

# SMDP: descuento variable por tiempo
mu, sigma = TIEMPO_ACCION[accion]   # ej: Norte = (2, 0.2)
tau = muestra_de_Normal(mu, sigma)  # ej: 1.87 segundos
desc = gamma ** tau                  # 0.97^1.87 = 0.944
\end{lstlisting}

\newpage

% ====================================================================
\section{Proyecto 3 — Q-Learning SMDP}
% ====================================================================

\subsection{Que es Q-Learning?}

\textbf{Q-Learning es aprender por prueba y error.}

Imagina un raton en un laberinto. Al principio no sabe nada — no sabe donde esta el queso, no sabe donde estan los muros. Solo puede probar caminos.

Despues de muchos intentos, empieza a aprender: ``si estoy en esta esquina y voy a la derecha, llego al queso''. Eso es Q-Learning.

\subsection{La tabla Q}

El robot tiene una tabla donde apunta notas:

\begin{lstlisting}
Q[(2,1), Norte]  = -3.2   (si estoy en (2,1) y voy Norte, fue malo)
Q[(2,1), Sur]    =  5.1   (si voy Sur, fue bueno!)
Q[(2,1), Este]   = -1.0   (neutro)
Q[(2,1), Oeste]  = -2.5   (malo)
\end{lstlisting}

Al principio \textbf{todos los valores son 0} (no sabe nada). Despues de muchos episodios, los valores se llenan y el robot sabe que hacer.

\subsection{La ecuacion de actualizacion}

\begin{equation}
Q(s,a) \leftarrow Q(s,a) + \alpha \left[ R + \gamma^\tau \cdot \max_{a'} Q(s',a') - Q(s,a) \right]
\end{equation}

\begin{definicion}[Desglose]
\begin{itemize}
    \item $Q(s,a)$ = mi nota actual para esta accion en este estado
    \item $\alpha = 0.10$ = cuanto ajusto (10\% = poco a poco)
    \item $R$ = lo que acaba de pasar ($-1$ o $+10$)
    \item $\gamma^\tau$ = descuento por tiempo (igual que SMDP)
    \item $\max Q(s')$ = lo mejor que puedo esperar del siguiente estado
    \item $Q(s,a)$ viejo = mi nota antes de esta experiencia
\end{itemize}
\end{definicion}

La diferencia entre ``lo que esperaba'' y ``lo que paso'' se llama \vocab{error temporal}. Si el error es grande, ajusto mucho. Si es chico, ajusto poco.

\newpage
\subsection{Epsilon-greedy (explorar vs explotar)}

El robot tiene un dilema:

\begin{itemize}
    \item \textbf{Explotar}: usar lo que ya sabe (ir por el mejor camino conocido)
    \item \textbf{Explorar}: probar algo nuevo (ir por un camino random)
\end{itemize}

\textbf{Epsilon} controla esto:

\begin{lstlisting}[language=Python]
def elegir(estado, epsilon):
    if random() < epsilon:          # ej: 30% de las veces
        return accion_aleatoria     # EXPLORAR: probar algo nuevo
    else:                           # ej: 70% de las veces
        return mejor_accion(estado) # EXPLOTAR
\end{lstlisting}

\begin{importante}
Al principio epsilon es alto (explorar mucho). Con el tiempo baja (explotar lo aprendido). El slider de epsilon te deja controlar esto en vivo.
\end{importante}

\subsection{Alpha (tasa de aprendizaje)}

$\alpha = 0.10$ significa que cada actualizacion solo cambia el valor en un 10\%:

\begin{lstlisting}
Q_viejo   = 5.0
nueva_info = 7.0
alpha     = 0.10

Q_nuevo = 5.0 + 0.10 * (7.0 - 5.0) = 5.0 + 0.2 = 5.2
\end{lstlisting}

Si $\alpha$ fuera 1.0, el nuevo valor seria 7.0 (borraria todo lo anterior). Con 0.10, cambia solo un poquito. Esto hace el aprendizaje \textbf{estable}.

\newpage
\subsection{El codigo del Q-Learning}

\begin{lstlisting}[language=Python]
class AgenteQLearning:
    def __init__(self):
        self.q = defaultdict(float)   # todo empieza en 0

    def entrenar_episodio(self, gamma, alpha, epsilon):
        s = get_random_start()        # empezar en random

        for _ in range(200):          # max 200 pasos
            a = self.elegir(s, epsilon)     # epsilon-greedy
            sig = self.mover(s, a)          # 90% exito
            r = -1 + (10 if sig == meta else 0)
            tau = self.tiempo_accion(a)     # muestrear tiempo
            desc = gamma ** tau             # descuento SMDP

            viejo = self.q[(s, a)]
            self.q[(s, a)] = viejo + alpha * (
                r + desc * self.max_q(sig) - viejo
            )

            s = sig
            if s == meta:
                break
\end{lstlisting}

Paso a paso:
\begin{enumerate}
    \item El robot empieza en una posicion random
    \item Elige una accion (epsilon-greedy)
    \item Se mueve (90\% exito, 10\% falla)
    \item Recibe recompensa ($-1$ por paso, $+10$ si llego a la meta)
    \item Mide cuanto tiempo tomo la accion
    \item Actualiza la tabla Q
    \item Repite hasta llegar a la meta o 200 pasos
    \item Repite todo desde paso 1 con otro inicio random
\end{enumerate}

Despues de \~{}1500 episodios, la tabla Q converge y el robot sabe que hacer en cada situacion.

\newpage

% ====================================================================
\section{Diferencias Clave entre los Proyectos}
% ====================================================================

\begin{center}
\begin{tabular}{>{\bfseries}l l l}
\toprule
\rowcolor{tabheader}
Aspecto & Proyecto 2 (MDP/SMDP) & Proyecto 3 (Q-Learning) \\
\midrule
Tipo & Basado en modelo & Libre de modelo \\
Conoce reglas? & SI & NO \\
Como aprende? & Calcula matematica & Prueba y error \\
Necesita modelo? & SI (P, R conocidos) & NO (solo experiencia) \\
Convergencia & Exacta, pocas iter & Aprox, miles episodios \\
Descuento & $\gamma$ o $\gamma^\tau$ & $\gamma^\tau$ siempre \\
Resultado & $V(s)$ y politica & Tabla $Q(s,a)$ y politica \\
\bottomrule
\end{tabular}
\end{center}

\begin{importante}
\textbf{Pregunta clave de defensa:} ``¿Cual es la diferencia entre MDP y Q-Learning?'' \\[6pt]
\textbf{Respuesta:} MDP necesita conocer las reglas del juego (las probabilidades y recompensas). Q-Learning no necesita nada — aprende jugando. Es como la diferencia entre estudiar el mapa antes de caminar vs aprender caminando.
\end{importante}

\newpage

% ====================================================================
\section{Cheat Sheet para la Defensa Oral}
% ====================================================================

\begin{pregunta}[1. Que es un MDP?]
Un modelo para tomar decisiones secuenciales bajo incertidumbre. 5 componentes: estados, acciones, probabilidades de transicion, recompensas y factor de descuento. La ecuacion de Bellman calcula el valor de cada estado y la mejor accion.
\end{pregunta}

\begin{pregunta}[2. Que es un SMDP y como se diferencia de un MDP?]
Un SMDP extiende el MDP al considerar que las acciones toman tiempo variable. Norte/Sur = N(2, 0.2), Este/Oeste = N(3, 0.3). El descuento cambia de $\gamma$ a $\gamma^\tau$, donde $\tau$ es el tiempo muestreado.
\end{pregunta}

\begin{pregunta}[3. Por que $\gamma^\tau$?]
Porque si una accion toma mas tiempo, la recompensa futura llega mas tarde y vale menos. $\tau$ grande = descuento grande = la recompensa futura vale menos.
\end{pregunta}

\begin{pregunta}[4. Que es Q-Learning?]
Algoritmo libre de modelo. El agente mantiene una tabla Q que estima el valor de cada par (estado, accion). Aprende por experiencia: elige acciones, observa recompensas, actualiza Q con: $Q(s,a) += \alpha \cdot [R + \gamma^\tau \cdot \max Q(s') - Q(s,a)]$.
\end{pregunta}

\begin{pregunta}[5. Que es alpha?]
Tasa de aprendizaje. $\alpha = 0.10$ = ajusta 10\% cada vez. Hace el aprendizaje gradual y estable — no borra todo lo que aprendio antes por una experiencia rara.
\end{pregunta}

\begin{pregunta}[6. Que es epsilon-greedy?]
Estrategia de exploracion vs explotacion. Con probabilidad $\epsilon$, el agente elige accion aleatoria (explorar). Con probabilidad $1-\epsilon$, elige la mejor conocida (explotar).
\end{pregunta}

\begin{pregunta}[7. Por que Q-Learning es ``libre de modelo''?]
No necesita conocer $P(s'|s,a)$ ni $R(s,a)$. Solo necesita interactuar con el entorno y observar que pasa. MDP necesita modelo completo; Q-Learning solo experiencia.
\end{pregunta}

\begin{pregunta}[8. Cuantos episodios para converger?]
$\sim$1500--2000 episodios. La recompensa promedio sube de $-17$ (al inicio) a $+1$ (cuando converge).
\end{pregunta}

\begin{pregunta}[9. Por que los valores SMDP $<$ MDP?]
Porque el tiempo extra en E/O (3s vs 1s) hace que $\gamma^\tau$ sea menor que $\gamma$. Descuento mas agresivo = valores menores.
\end{pregunta}

\begin{pregunta}[10. Por que inicio aleatorio?]
Para que funcione desde cualquier estado. Q-Learning necesita explorar desde distintos puntos para llenar bien la tabla Q.
\end{pregunta}

\newpage

% ====================================================================
\section{Rubricas de Evaluacion}
% ====================================================================

\subsection{Proyecto 2 (100 puntos)}

\begin{center}
\begin{tabular}{l c l}
\toprule
\rowcolor{tabheader}
Criterio & Pts & Que necesitas \\
\midrule
Simulacion grafica & 20 & Pygame, mapa 2D, robot, meta, muros \\
MDP clasico & 30 & Value Iteration, $\gamma=0.97$, P=90\%, 4 acciones \\
SMDP & 30 & $\gamma^\tau$, N(2,0.2), N(3,0.3), 90\% exito \\
Defensa oral & 20 & Explicar diferencia MDP/SMDP, dominar codigo \\
\bottomrule
\end{tabular}
\end{center}

\subsection{Proyecto 3 (100 puntos)}

\begin{center}
\begin{tabular}{l c l}
\toprule
\rowcolor{tabheader}
Criterio & Pts & Que necesitas \\
\midrule
Simulacion grafica & 20 & Pygame, mapa 2D, convergencia visible \\
Q-Learning SMDP & 40 & Tabla Q, $\alpha=0.10$, $\gamma^\tau$, epsilon-greedy \\
Modelado estocastico & 20 & 90\% exito, N(2,0.2), N(3,0.3) \\
Defensa oral & 20 & Explicar Q-Learning vs MDP, dominar codigo \\
\bottomrule
\end{tabular}
\end{center}

\newpage

% ====================================================================
\section{Parametros y Controles}
% ====================================================================

\subsection{Sliders (ajustables en vivo)}

\begin{center}
\begin{tabular}{l c c l}
\toprule
\rowcolor{tabheader}
Slider & Proy 2 & Proy 3 & Que hace \\
\midrule
Gamma & Si & Si & Factor de descuento (0.5 a 0.99) \\
Alpha & No & Si & Tasa de aprendizaje (0.01 a 0.50) \\
Epsilon & No & Si & Exploracion (0.01 a 1.0) \\
Velocidad & Si & Si & Ms entre pasos (100 a 800) \\
\bottomrule
\end{tabular}
\end{center}

\subsection{Controles Proyecto 2}

\begin{tabular}{ll}
\textbf{1} o boton MDP & Cambia a modo MDP clasico \\
\textbf{2} o boton SMDP & Cambia a modo SMDP \\
\textbf{ESPACIO} & Pausar / Reanudar \\
\textbf{R} o boton Reset & Reiniciar \\
\end{tabular}

\subsection{Controles Proyecto 3}

\begin{tabular}{ll}
\textbf{ENTER} o boton Entrenar & 500 episodios \\
\textbf{T} o boton Auto & Hasta convergencia \\
\textbf{A} o boton Animar & Ver robot navegar \\
\textbf{R} o boton Reset & Reiniciar todo \\
\end{tabular}

\subsection{Como correr}

\begin{lstlisting}
pip install pygame numpy
cd C:\Projects\robotica
python proyecto2.py
python proyecto3.py
\end{lstlisting}

% ====================================================================
% APENDICE: Formulas rapidas
% ====================================================================

\newpage
\section{Apéndice: Formulas Rapidas}

\begin{center}
\begin{tabular}{>{\bfseries}l l}
\toprule
\rowcolor{tabheader}
Concepto & Formula \\
\midrule
Bellman (MDP) & $V(s) = \max_a [R + \gamma \cdot V(s')]$ \\
Bellman (SMDP) & $V(s) = \max_a [R + \mathbb{E}[\gamma^\tau] \cdot V(s')]$ \\
Q-Learning & $Q(s,a) += \alpha [R + \gamma^\tau \max Q(s') - Q(s,a)]$ \\
Descuento temporal & $\gamma^\tau = 0.97^\tau$ \\
Transicion (90\%) & $0.9 \cdot V(s') + 0.1 \cdot V(s)$ \\
Recompensa & $R = -1$ (paso), $R = +10$ (meta) \\
\bottomrule
\end{tabular}
\end{center}

\end{document}
